\section{Related Work}
\label{sec:related}

\subsection{Data-Driven Fault Diagnosis in Nuclear Plants}

Machine-learning approaches to nuclear fault diagnosis have evolved from classical pattern-recognition methods to deep sequence models. Early work used feed-forward networks and support vector machines on static snapshots of sensor data, and was reviewed comprehensively by Santhosh~et~al.~\citep{santhosh2019survey}. The release of the open-source Nuclear Power Plant Accident Dataset~\citep{qi2022nppad} enabled the systematic benchmarking of modern sequence architectures. Subsequent studies applied convolutional and recurrent networks to the 18-class accident identification problem, reporting high in-distribution accuracy for common transients but persistent difficulties with rare or physically similar accident types~\citep{wang2023deep,lee2023lstm}.

More recent work has turned to explainability. A January 2026 study by \placeholder{Author et al.}~\citep{xgboost2026shap} combined SHAP-based feature attribution with an incremental XGBoost classifier on a \nppad-derived feature set, using SHAP primarily for feature selection and incremental model updating. Our work differs in two ways. First, we operate on raw multivariate windows rather than hand-engineered feature summaries, retaining the temporal structure that downstream prognosis requires. Second, we use SHAP not as a feature-selection heuristic but as a physics-grounded root-cause-attribution tool---validating that the classifier's attributions align with the known mechanistic signatures of each accident type---and we fold the attribution analysis into a broader pipeline rather than treating it as a standalone result.

\subsection{Physics-Informed Machine Learning for Nuclear Systems}

Physics-informed neural networks (PINNs) augment data-driven objectives with residuals from governing partial or ordinary differential equations~\citep{raissi2019pinn,karniadakis2021pinn}. In the nuclear domain, PINNs have been applied to reactor transient prediction~\citep{zhang2023pinn}, autonomous core calibration~\citep{song2023calibration}, and reduced-order modeling within digital-twin workflows~\citep{gong2022rom}. These studies demonstrate that embedding reactor physics (point kinetics, two-phase flow closures, or full neutron transport) improves generalization to regimes with limited training data.

Our physics-informed surrogate follows this tradition but differs in its embedding scope and pipeline role. We embed point kinetics and steady-state energy conservation constraints---rather than a full neutron-transport operator---because our target is a \emph{plant-level} surrogate covering 96 coupled sensor channels rather than a detailed core model. More importantly, our surrogate is designed from the outset as the third stage of a diagnostic pipeline; its outputs are used as consistency checks on the classifier, not only as physically plausible forecasts in their own right.

\subsection{Digital Twins and Integrated Pipelines for NPPs}

The digital-twin paradigm for nuclear plants has been the subject of several recent review articles~\citep{ayodeji2022dtwin,digitaltwin2025review}, which converge on a monitoring--diagnosis--prognosis architecture as the natural software organization for safety-relevant decision support. These reviews emphasize, however, that few published works actually implement the full chain: most demonstrations terminate at classification, and the prognosis stage is often treated as future work. Where prognosis is addressed, it is typically a purely data-driven forecaster without explicit physics constraints, which makes its predictions difficult to trust for out-of-distribution transients or unseen accident severities.

Our contribution is a concrete implementation of the monitoring--diagnosis--prognosis chain on a public dataset, with explicit physics constraints in the prognosis stage and an attribution layer in the diagnosis stage that ties model decisions back to mechanistic knowledge. This makes the system interpretable at the level of individual predictions and auditable at the level of physical consistency, both of which are prerequisites for regulated deployment.
